%%% BEGIN fig:twist3-ssa
\begin{figure}[htbp]
\centering
\begin{tikzpicture}[fd,
  t3 hadron/.style={draw=black!70,line width=1.5pt,line cap=butt},
  t3 pole/.style={draw=red!70!black,line width=1.6pt,fd arrow=#1},
  t3 pole/.default=0.5,
  t3 lab/.style={font=\small},
  t3 note/.style={font=\footnotesize}]
  \node[jet blob,ellipse,minimum width=5.6cm,minimum height=1.1cm,inner sep=0pt] (A) at (0,-2.2) {};
  \node[jetB blob,ellipse,minimum width=5.6cm,minimum height=1.1cm,inner sep=0pt] (B) at (0,2.3) {};
  \node[hard blob,circle,minimum size=0.84cm,inner sep=0pt] (HL) at (-2.55,0) {};
  \node[hard blob,circle,minimum size=0.84cm,inner sep=0pt] (HR) at (2.55,0) {};
  \coordinate (G) at (-0.7,0);
  % hadrons
  \draw[t3 hadron] (-4.75,-2.2) -- (A.west);  \draw[t3 hadron] (A.east) -- (4.75,-2.2);
  \draw[t3 hadron] (-4.75, 2.3) -- (B.west);  \draw[t3 hadron] (B.east) -- (4.75, 2.3);
  % partons from B (gluons) and from A (quarks x1, x2)
  \draw[gluon] (-1.75,1.87) -- (HL.105);   \draw[gluon] (HR.75) -- (1.75,1.87);
  \draw[quark] (-1.75,-1.77) -- (HL.285);   \draw[quark] (HR.255) -- (1.75,-1.77);
  % soft gluon of the twist-3 correlator, attached to the final-state quark c
  \draw[gluon] (-0.7,-1.67) -- (G);
  % final-state quark c (pole propagator in red) and the observed photon
  \draw[t3 pole=0.55] (HL.east) -- (G);
  \draw[quark=0.5] (G) -- (HR.west);
  \draw[photon] (HL.55) .. controls (-1.3,1.2) and (1.3,1.2) .. (HR.125);
  \node[vertex] at (G) {};
  \draw[cut] (0,-2.95) -- (0,3.05);
  % labels
  \node[t3 lab] at (HL) {$H$};  \node[t3 lab] at (HR) {$H^\dagger$};
  \node[t3 lab] at (-1.3,-2.25) {$A^{\uparrow}$};  \node[t3 lab] at (-1.3,2.3) {$B$};
  \node[t3 lab,below=1pt] at (-4.3,-2.25) {$P_A,S_{\mathrm T}$};
  \node[t3 lab,above=1pt] at (-4.3,2.35) {$P_B$};
  \path[mom={$x_1P_A$},mom sep=0.27cm,mom length=0.45cm] (-1.75,-1.77) -- (HL.285);
  \path[mom={$x_2P_A$},mom sep=0.27cm,mom length=0.45cm] (HR.255) -- (1.75,-1.77);
  \path[mom={$x_gP_A$},mom sep=0.27cm,mom length=0.45cm,mom pos=0.45] (-0.7,-1.67) -- (G);
  \path[mom={$p_c$},mom sep=0.25cm,mom length=0.4cm] (G) -- (0,0);
  \node[t3 lab] at (0.62,1.22) {$\gamma$};
  \node[t3 note,text=red!65!black] at (-1.42,0.25) {极点};
  \node[t3 note,anchor=north] at (0,-2.95) {$T_F(x_1,x_2)$：夸克--胶子--夸克关联（扭度 3）};
\end{tikzpicture}
\caption{Qiu--Sterman 软胶子极点机制（以 $A^{\uparrow}+B\to\gamma+X$ 中的 Compton 子过程 $qg\to\gamma q$ 为例；竖直虚线为末态割线）。
横向极化强子 $A^{\uparrow}$ 的扭度 3 关联函数 $T_F(x_1,x_2)$ 提供割线两侧的夸克线（$x_1P_A$、$x_2P_A$）与一条胶子线（$x_gP_A$，$x_g\equiv x_2-x_1$），后者与末态夸克 $c$ 相连。
夹在硬部分与该胶子之间的传播子（红色）携带 $p_c-(x_2-x_1)P_A$，其分母 $\big(p_c-(x_2-x_1)P_A\big)^2+\ii0=-2p_c\cdot P_A\,(x_2-x_1-\ii0)$ 在 $x_2\to x_1$ 时为零：
$\frac{1}{x_2-x_1-\ii0}$ 的 $\ii\pi\delta(x_2-x_1)$ 部分提供单自旋不对称所需的相位，使结果正比于 $T_F(x,x)$（“软胶子极点”）。
若胶子与初态部分子相连，$\ii0$ 的符号相反，色因子也不同——这与 Sivers 函数的过程依赖性相对应（参见\cref{fig:sidis-dy-links}）。}
\label{fig:twist3-ssa}
\end{figure}
%%% END fig:twist3-ssa
